\documentclass{article}
\usepackage{amsmath,amssymb,amsthm,algorithm,algorithmic}
\usepackage{hyperref}
\newtheorem{theorem}{Theorem}
\newtheorem{lemma}{Lemma}
\newtheorem{assumption}{Assumption}
\newcommand{\prox}{\operatorname{prox}}
\newcommand{\grad}{\nabla}
\begin{document}

\section*{Proximal Gradient Method (PGM)}

\section*{Mathematical Formulation}
Let $f:\mathbb{R}^{d}\to\mathbb{R}$ be convex and $L$‑smooth, i.e.
\[
\|\grad f(x)-\grad f(y)\|\le L\|x-y\|,\qquad\forall x,y\in\mathbb{R}^{d}.
\]
Let $g:\mathbb{R}^{d}\to\mathbb{R}\cup\{+\infty\}$ be a proper, closed, convex (possibly non‑smooth) function.  
Define the composite objective
\[
F(x):=f(x)+g(x).
\]
For a stepsize $\eta>0$ the proximal gradient iteration is
\begin{equation}
\label{eq:pgm}
x_{k+1}= \prox_{\eta g}\bigl(x_{k}-\eta \grad f(x_{k})\bigr),\qquad k=0,1,2,\dots
\end{equation}
where the proximal operator of $g$ is
\[
\prox_{\eta g}(u):=\arg\min_{z\in\mathbb{R}^{d}}\Bigl\{ g(z)+\frac{1}{2\eta}\|z-u\|^{2}\Bigr\}.
\]

\section*{Pseudocode}
\begin{algorithm}[H]
\caption{Proximal Gradient Method}
\begin{algorithmic}[1]
\STATE \textbf{Input:} stepsize $\eta\in(0,1/L]$, initial point $x_{0}\in\mathbb{R}^{d}$, max iterations $K$.
\FOR{$k=0$ \textbf{to} $K-1$}
\STATE Compute gradient $g_{k}\gets \grad f(x_{k})$.
\STATE Take a gradient step $y_{k}\gets x_{k}-\eta g_{k}$.
\STATE Apply proximal map $x_{k+1}\gets \prox_{\eta g}(y_{k})$.
\ENDFOR
\STATE \textbf{Output:} $x_{K}$ (or the averaged iterate $\bar x_{K}=\frac1K\sum_{k=1}^{K}x_{k}$).
\end{algorithmic}
\end{algorithm}

\section*{Dry Run Example}
Consider the scalar problem
\[
f(w)=(w-3)^{2},\qquad g\equiv 0,\qquad \eta=0.1,\qquad w_{0}=0.
\]
Since $g\equiv0$, $\prox_{\eta g}(u)=u$.

\begin{center}
\begin{tabular}{c|c|c|c}
Iteration $k$ & $w_{k}$ & $\grad f(w_{k})=2(w_{k}-3)$ & $w_{k+1}=w_{k}-\eta\grad f(w_{k})$ \\ \hline
0 & $0$ & $-6$ & $0-0.1(-6)=0.6$ \\ 
1 & $0.6$ & $2(0.6-3)=-4.8$ & $0.6-0.1(-4.8)=1.08$ \\ 
2 & $1.08$ & $2(1.08-3)=-3.84$ & $1.08-0.1(-3.84)=1.464$ \\ 
3 & $1.464$ & $2(1.464-3)=-3.072$ & $1.464-0.1(-3.072)=1.7712$
\end{tabular}
\end{center}
The objective values $f(w_{k})$ decrease:
\[
f(0)=9,\; f(0.6)=5.76,\; f(1.08)=3.6864,\; f(1.464)=2.2143.
\]

\newpage
\section*{Assumptions}
\begin{assumption}[Smoothness of $f$]\label{ass:smooth}
$f$ is convex and $L$‑smooth, i.e. for all $x,y$,
\[
f(y)\le f(x)+\langle \grad f(x),y-x\rangle+\frac{L}{2}\|y-x\|^{2}.
\]
\end{assumption}
\begin{assumption}[Convexity of $g$]\label{ass:gconv}
$g$ is proper, closed and convex; its proximal operator $\prox_{\eta g}$ is well defined for any $\eta>0$.
\end{assumption}
\begin{assumption}[Stepsize]\label{ass:stepsize}
The stepsize satisfies $0<\eta\le \frac{1}{L}$.
\end{assumption}
\begin{assumption}[Existence of a minimizer]\label{ass:opt}
There exists at least one optimal solution $x^{\star}\in\arg\min_{x}F(x)$ and $F(x^{\star})<+\infty$.
\end{assumption}

\section*{Main Convergence Theorem}
\begin{theorem}[Sub‑linear rate of PGM]\label{thm:sublinear}
Under Assumptions \ref{ass:smooth}–\ref{ass:opt}, the iterates generated by \eqref{eq:pgm} satisfy for all $k\ge 1$,
\[
F(x_{k})-F(x^{\star})\le \frac{\|x_{0}-x^{\star}\|^{2}}{2\eta k}.
\]
Equivalently, the averaged iterate $\bar x_{k}:=\frac1k\sum_{t=0}^{k-1}x_{t}$ fulfills
\[
F(\bar x_{k})-F(x^{\star})\le \frac{\|x_{0}-x^{\star}\|^{2}}{2\eta k}.
\]
\end{theorem}

\section*{Theoretical Properties}
\begin{itemize}
\item \textbf{Stability:} Because $\prox_{\eta g}$ is non‑expansive,
\[
\|\prox_{\eta g}(u)-\prox_{\eta g}(v)\|\le\|u-v\|,\qquad\forall u,v,
\]
the iteration \eqref{eq:pgm} is a non‑expansive map composed with a gradient step that is a contraction for $\eta\le 1/L$.
\item \textbf{Hyperparameter sensitivity:} The bound in Theorem \ref{thm:sublinear} scales linearly with $1/\eta$. Choosing $\eta$ as large as allowed ($\eta=1/L$) yields the best constant.
\item \textbf{Comparison to Gradient Descent:} When $g\equiv0$, $\prox_{\eta g}$ reduces to the identity and the theorem recovers the classical $O(1/k)$ rate for gradient descent on convex smooth functions.
\item \textbf{Extension to Strong Convexity:} If $F$ is $\mu$‑strongly convex, the same proof (with Lemma \ref{lem:strong}) yields a linear rate $O((1-\eta\mu)^{k})$.
\end{itemize}

\section*{Discussion}
The proof hinges on two elementary facts:
\begin{enumerate}
\item The \emph{descent lemma} for $L$‑smooth $f$ (Lemma \ref{lem:descent}).
\item The \emph{optimality condition} of the proximal map, which together with its non‑expansiveness produces a one‑step inequality that telescopes.
\end{enumerate}
No stochasticity is introduced, therefore expectations are not needed; the derivation is deterministic and holds for any initialization.

\newpage
\section*{Preliminaries (Definitions and Lemmas)}
\begin{lemma}[Descent Lemma]\label{lem:descent}
If $f$ is $L$‑smooth, then for any $x,y$,
\[
f(y)\le f(x)+\langle \grad f(x),y-x\rangle+\frac{L}{2}\|y-x\|^{2}.
\]
\end{lemma}
\begin{lemma}[Proximal optimality]\label{lem:prox_opt}
For any $u\in\mathbb{R}^{d}$ and $\eta>0$, $z^{+}:=\prox_{\eta g}(u)$ satisfies
\[
0\in \partial g(z^{+})+\frac{1}{\eta}(z^{+}-u),
\]
i.e.
\[
\langle z^{+}-u,\,z-z^{+}\rangle +\eta\bigl(g(z)-g(z^{+})\bigr)\ge 0,\qquad \forall z\in\mathbb{R}^{d}.
\]
\end{lemma}
\begin{lemma}[Non‑expansiveness of the proximal operator]\label{lem:nonexp}
For any $u,v\in\mathbb{R}^{d}$,
\[
\|\prox_{\eta g}(u)-\prox_{\eta g}(v)\|\le\|u-v\|.
\]
\end{lemma}
\begin{lemma}[Strong convexity (optional)]\label{lem:strong}
If $F$ is $\mu$‑strongly convex, then
\[
F(x)-F(x^{\star})\ge \frac{\mu}{2}\|x-x^{\star}\|^{2}.
\]
\end{lemma}

\section*{Main Proof (Step‑by‑Step Derivation)}
Let $x^{\star}$ be a minimizer of $F$. Define the gradient step
\[
y_{k}=x_{k}-\eta\grad f(x_{k}),\qquad x_{k+1}= \prox_{\eta g}(y_{k}).
\]

\noindent\textbf{[Step 1]} Apply Lemma \ref{lem:prox_opt} with $u=y_{k}$, $z^{+}=x_{k+1}$ and arbitrary $z=x^{\star}$:
\begin{equation}
\label{eq:opt_cond}
\langle x_{k+1}-y_{k},\,x^{\star}-x_{k+1}\rangle+\eta\bigl(g(x^{\star})-g(x_{k+1})\bigr)\ge 0.
\end{equation}

\noindent\textbf{[Step 2]} Expand the inner product:
\[
\langle x_{k+1}-y_{k},\,x^{\star}-x_{k+1}\rangle
= \langle x_{k+1}-y_{k},x^{\star}-y_{k}\rangle-\|x_{k+1}-y_{k}\|^{2}.
\]
Insert this into \eqref{eq:opt_cond} and rearrange:
\begin{equation}
\label{eq:basic}
\|x_{k+1}-y_{k}\|^{2}\le \langle x_{k+1}-y_{k},x^{\star}-y_{k}\rangle+\eta\bigl(g(x^{\star})-g(x_{k+1})\bigr).
\end{equation}

\noindent\textbf{[Step 3]} Replace $y_{k}=x_{k}-\eta\grad f(x_{k})$; then $x_{k+1}-y_{k}=x_{k+1}-x_{k}+ \eta\grad f(x_{k})$. Hence
\[
\|x_{k+1}-y_{k}\|^{2}= \|x_{k+1}-x_{k}\|^{2}+2\eta\langle\grad f(x_{k}),x_{k+1}-x_{k}\rangle+\eta^{2}\|\grad f(x_{k})\|^{2}.
\]

\noindent\textbf{[Step 4]} Substitute the expression above into the left‑hand side of \eqref{eq:basic} and move all terms to one side:
\begin{align}
\|x_{k+1}-x_{k}\|^{2}&+2\eta\langle\grad f(x_{k}),x_{k+1}-x_{k}\rangle+\eta^{2}\|\grad f(x_{k})\|^{2} \nonumber\\
&\le \langle x_{k+1}-x_{k}+\eta\grad f(x_{k}),\,x^{\star}-x_{k}+\eta\grad f(x_{k})\rangle \nonumber\\
&\qquad +\eta\bigl(g(x^{\star})-g(x_{k+1})\bigr). \label{eq:after_sub}
\end{align}

\noindent\textbf{[Step 5]} Expand the right‑hand side of \eqref{eq:after_sub}:
\begin{align}
\text{RHS}= &\ \langle x_{k+1}-x_{k},x^{\star}-x_{k}\rangle
+\eta\langle x_{k+1}-x_{k},\grad f(x_{k})\rangle \nonumber\\
&+\eta\langle \grad f(x_{k}),x^{\star}-x_{k}\rangle
+\eta^{2}\|\grad f(x_{k})\|^{2}
+\eta\bigl(g(x^{\star})-g(x_{k+1})\bigr). \label{eq:rhs_expand}
\end{align}

\noindent\textbf{[Step 6]} Cancel the identical $\eta^{2}\|\grad f(x_{k})\|^{2}$ terms on both sides of \eqref{eq:after_sub}. After cancellation, bring the remaining terms to the left:
\begin{align}
\|x_{k+1}-x_{k}\|^{2}
&+2\eta\langle\grad f(x_{k}),x_{k+1}-x_{k}\rangle
-\eta\langle x_{k+1}-x_{k},\grad f(x_{k})\rangle \nonumber\\
&\le \langle x_{k+1}-x_{k},x^{\star}-x_{k}\rangle
+\eta\langle \grad f(x_{k}),x^{\star}-x_{k}\rangle
+\eta\bigl(g(x^{\star})-g(x_{k+1})\bigr). \label{eq:mid}
\end{align}
Observe that $2\eta\langle\grad f(x_{k}),x_{k+1}-x_{k}\rangle-\eta\langle x_{k+1}-x_{k},\grad f(x_{k})\rangle = \eta\langle\grad f(x_{k}),x_{k+1}-x_{k}\rangle$.

Thus \eqref{eq:mid} simplifies to
\begin{equation}
\label{eq:simplified}
\|x_{k+1}-x_{k}\|^{2}
+\eta\langle\grad f(x_{k}),x_{k+1}-x_{k}\rangle
\le \langle x_{k+1}-x_{k},x^{\star}-x_{k}\rangle
+\eta\langle \grad f(x_{k}),x^{\star}-x_{k}\rangle
+\eta\bigl(g(x^{\star})-g(x_{k+1})\bigr).
\end{equation}

\noindent\textbf{[Step 7]} Rearrange the inner products:
\[
\langle x_{k+1}-x_{k},x^{\star}-x_{k}\rangle
= \frac12\Bigl(\|x^{\star}-x_{k}\|^{2}-\|x^{\star}-x_{k+1}\|^{2}-\|x_{k+1}-x_{k}\|^{2}\Bigr).
\]
Insert this identity into \eqref{eq:simplified} and cancel the $\|x_{k+1}-x_{k}\|^{2}$ terms:
\begin{align}
\eta\langle\grad f(x_{k}),x_{k+1}-x_{k}\rangle
&\le \frac12\bigl(\|x^{\star}-x_{k}\|^{2}-\|x^{\star}-x_{k+1}\|^{2}\bigr) \nonumber\\
&\quad +\eta\langle \grad f(x_{k}),x^{\star}-x_{k}\rangle
+\eta\bigl(g(x^{\star})-g(x_{k+1})\bigr). \label{eq:after_id}
\end{align}

\noindent\textbf{[Step 8]} Move the term $\eta\langle\grad f(x_{k}),x_{k+1}-x_{k}\rangle$ to the right‑hand side and combine the two gradient inner products:
\begin{equation}
\label{eq:combined}
\frac12\bigl(\|x^{\star}-x_{k}\|^{2}-\|x^{\star}-x_{k+1}\|^{2}\bigr)
\ge \eta\bigl(f(x_{k+1})-f(x_{k})\bigr)
+\eta\bigl(g(x^{\star})-g(x_{k+1})\bigr).
\end{equation}
The inequality uses the \emph{descent lemma} (Lemma \ref{lem:descent}) with $y=x_{k+1}$:
\[
f(x_{k+1})\le f(x_{k})+\langle\grad f(x_{k}),x_{k+1}-x_{k}\rangle+\frac{L}{2}\|x_{k+1}-x_{k}\|^{2},
\]
and the stepsize condition $\eta\le 1/L$ to drop the quadratic term, i.e.
\[
\eta\langle\grad f(x_{k}),x_{k+1}-x_{k}\rangle\ge \eta\bigl(f(x_{k+1})-f(x_{k})\bigr).
\]

\noindent\textbf{[Step 9]} Rearrange \eqref{eq:combined} to isolate the objective gap:
\[
\eta\bigl(F(x_{k+1})-F(x^{\star})\bigr)
\le \frac12\bigl(\|x^{\star}-x_{k}\|^{2}-\|x^{\star}-x_{k+1}\|^{2}\bigr).
\]

\noindent\textbf{[Step 10]} Divide by $\eta$ and sum the inequality from $k=0$ to $k=K-1$:
\[
\sum_{k=0}^{K-1}\bigl(F(x_{k+1})-F(x^{\star})\bigr)
\le \frac{1}{2\eta}\bigl(\|x^{\star}-x_{0}\|^{2}-\|x^{\star}-x_{K}\|^{2}\bigr)
\le \frac{\|x^{\star}-x_{0}\|^{2}}{2\eta}.
\]

\noindent\textbf{[Step 11]} Since the left‑hand side is a sum of $K$ non‑negative terms, the smallest term is bounded by the average:
\[
\min_{1\le k\le K}\bigl(F(x_{k})-F(x^{\star})\bigr)
\le \frac1K\sum_{k=1}^{K}\bigl(F(x_{k})-F(x^{\star})\bigr)
\le \frac{\|x^{\star}-x_{0}\|^{2}}{2\eta K}.
\]
Because $F$ is convex, the same bound holds for the averaged iterate $\bar x_{K}$, completing the proof of Theorem \ref{thm:sublinear}.

\section*{Rate Derivation (With Constants)}
From the final bound we obtain the explicit rate
\[
F(x_{k})-F(x^{\star})\le \frac{R^{2}}{2\eta k},
\qquad R:=\|x_{0}-x^{\star}\|.
\]
If we set the maximal admissible stepsize $\eta=1/L$, the constant becomes $L R^{2}/2$:
\[
F(x_{k})-F(x^{\star})\le \frac{L\,\|x_{0}-x^{\star}\|^{2}}{2k}.
\]

\section*{Special Cases}
\begin{itemize}
\item \textbf{Pure gradient descent ($g\equiv0$).}  The proximal map is the identity, and the bound reduces to the classical $O(1/k)$ rate for smooth convex optimization.
\item \textbf{Indicator function $g=\iota_{C}$ of a closed convex set $C$.}  $\prox_{\eta g}$ becomes the Euclidean projection onto $C$, yielding the projected gradient method with the same $O(1/k)$ guarantee.
\item \textbf{Strongly convex $F$.}  Using Lemma \ref{lem:strong} together with the inequality in Step 9 gives
\[
\|x_{k+1}-x^{\star}\|^{2}\le (1-\eta\mu)\|x_{k}-x^{\star}\|^{2},
\]
hence a linear convergence rate $F(x_{k})-F(x^{\star})\le (1-\eta\mu)^{k}\frac{L}{2}\|x_{0}-x^{\star}\|^{2}$.
\end{itemize}

\end{document}